\section{Conclusion}
\label{sec:conclusion}

We studied how supplementary evidence should enter a frozen concept-grounded segmentation model whose Presence, Semantic, and Instance readouts serve different functions. The central finding is that a shared concept supplies a stable category anchor but does not impose a shared evidence requirement. JRV preserves the official ontology and the native $p\max(S,I)$ path, specializes supplementary text by readout role, and composes aligned Local and Context observations after complete grounding. The eight-benchmark evaluation supports this organization: average mIoU rises from 53.56 to 57.49, every dataset improves over the matched baseline, and the fixed Role--View configuration has positive interaction on six datasets. The category and pixel analyses show that the gain reflects corrected assignments rather than a single dominant class.

The scope is equally part of the result. Dataset-level profiles are selected offline rather than predicted per image, the view operator varies with data distribution, and iSAID and UDD5 exhibit overlapping rather than complementary corrections. Cityscapes selects a Global endpoint, while the same-image Potsdam/Vaihingen control does not support an additional gain from the tested 256/512 crops. These observations position JRV as evidence organization for a frozen concept-grounded model, not as a universal claim that more prompts or larger fields improve segmentation. A natural next step is to replace offline profile selection with a validation-driven or lightweight adaptive policy while preserving the anchor and full-grounding constraints identified here.
